%------------------------------------------------------------------------------%
\begin{exercises}

%------------------------------------------------------------------------------%
\begin{exercise}
  Calcule as Integrais Indefinidas.
  \begin{mathlist}
    \item \int x^4- \frac{1}{2}x^3+\frac{1}{4}x-2 \,dx
    \item \int \sqrt{x^3}+ \sqrt[3]{x^2} \,dx
    \item \int \sqrt{x^2}+ x^{-2} \,dx
    \item \int (u+4)(2u+1) \,du
    \item \int v\left(v^2+2\right)^2 \,dv
    \item \int \frac{x^3-2 \sqrt{x}}{x} \,dx
    \item \int x^2 + 1 + \frac{1}{1+x^2} - 100\cos(x) \,dx
    \item \int \sin(2x)-2^x \,dx
    \item \int \theta - \csc(\theta)\cot(\theta) \,d\theta
    \item \int 1 + 10\tan(x)\sec(x) \,dx
    \item \int \csc^2(t) - 2e^t - e^{-t} \,dt
    \item \int \sec(t) \big( \sec(t) + \tan(t) \big) \,dt
    \item \int \frac{x^2-x+1}{x^3} dx
    \item \int \frac{\sin(2x)}{\sin(x)} \,dx
  \end{mathlist}
  \begin{answer}
    \begin{mathlist}
      \item \frac{x^5}{5}-\frac{x^4}{8}+\frac{x^2}{8}-2x+c
      \item \frac{2}{5}x^{\sfrac{5}{2}}+\frac{3}{5}x^{\sfrac{5}{3}}+C
      \item \frac{x^3}{3}-\frac{1}{x}+c
      \item \frac{2}{3}u^3+ \frac{9}{2}u^2+4u+C
      \item \frac{v^6}{6} +v^4+2v^2+C
      \item \frac{x^3}{3} - 4\sqrt{x}+c
      \item \frac{x^3}{3}+x+ \arctan(x)+100\sin(x)+C
      \item -\frac{\cos(2x)}{2}-\frac{2^x}{\ln(2)}+C
      \item \frac{\theta ^2}{2}+ \csc(x)+C
      \item x+10 \sec(x)+C
      \item -\cot(t)-2e^t+e^{-t}+C
      \item \tan(t)+\sec(t)+C
      \item \ln\abs{x} + x^{-1} - \frac{x^{-2}}{2} + C
      \item 2\cos(x)+C
    \end{mathlist}
  \end{answer}
\end{exercise}

%------------------------------------------------------------------------------%
\begin{exercise}
  Uma partícula se move com aceleração dada por 
  \[
    a(t)= 10 \sin(t)+ 3 \cos(t) \,.
  \]
  Encontre a equação do espaço sabendo que $s(0)=0$ e $s(2\pi)=12$.
  \begin{answer}
    \(
      s(t) =
      - 10\sin(t)
      - 3\cos(t)
      + \frac{6}{\pi} t
      + 3\)
  \end{answer}
\end{exercise}

%------------------------------------------------------------------------------%
\begin{exercise}
  Mostre que, para um movimento em uma reta com aceleração constante $a$,
  velocidade inicial $v_0$ e deslocamento inicial $s_0$, 
  o deslocamento no tempo $t$ é
  \[
    s = \frac{1}{2}at^2 + v_0t + s_0 \,.
  \]
\end{exercise}

%------------------------------------------------------------------------------%
%\begin{exercise}
%\begin{answer}
%\end{answer}
%\end{exercise}

%------------------------------------------------------------------------------%
%\begin{exercise}
%\begin{mathlist}[2]
%  \item
%  \item
%  \item
%\end{mathlist}
%\begin{answer}
%  \begin{alphalist}[3]
  %    \item
  %    \item
  %    \item
  %  \end{alphalist}
%\end{answer}
%\end{exercise}

\end{exercises}
%------------------------------------------------------------------------------%
